\documentclass[11pt]{article}
% Standalone preamble for web-only supplementary notes.
\usepackage[T1]{fontenc}
\usepackage{amsmath,amssymb,amsthm}
\usepackage{hyperref}
\usepackage{cleveref}
\DeclareMathOperator{\Map}{Map}
\DeclareMathOperator{\Fun}{Fun}
\newcommand{\iso}{\xrightarrow{\sim}}
\newtheorem{lemma}{Supplementary Lemma}
\newtheorem{corollary}[lemma]{Supplementary Corollary}
\crefname{lemma}{Supplementary Lemma}{Supplementary Lemmas}
\Crefname{lemma}{Supplementary Lemma}{Supplementary Lemmas}
\crefname{corollary}{Supplementary Corollary}{Supplementary Corollaries}
\Crefname{corollary}{Supplementary Corollary}{Supplementary Corollaries}

\begin{document}
\section*{Base change of triangles}

Base change must carry the specified triangle of a functor over the base,
not just its underlying functor. We spell out the calculation supporting
\href{exponentiable-functors.html}{base change of relative functors}.
Use the Chapter 1 structures and the conventions of
\href{supplementary-triangles-over-a-base.html}{Triangles over a base}.

Fix $p\colon S\to T$ and categories $f\colon C\to T$ and $g\colon D\to T$.
Write $r\colon C\times_T S\to C$, $q\colon C\times_T S\to S$ and
$\tau\colon fr=pq$ for the chosen pullback data. If $u=(h,\theta)$ is a
functor from $f$ to $g$ over $T$, form the cone with legs $hr,q$ and matching
\[
 g(hr)\simeq(gh)r\xrightarrow{\theta r}fr\xrightarrow{\tau}pq.
\]
Its lift defines $p^*u$. The second projection computation is its specified
triangle over $S$.

\begin{lemma}
\label[lemma]{lem:Supplementary_Base_Change_Comparison}
A comparison $u\simeq_T v$ induces a comparison $p^*u\simeq_S p^*v$.
\end{lemma}
\begin{proof}[Proof idea]
Write $v=(k,\psi)$ and let $\alpha\colon h=k$ be the given comparison.
Restrict $\psi\cdot(g\alpha)=\theta$ along $r$, and compare the two
associations of the resulting composites. Pasting with $\tau$ gives the
matching equation for the two cones. This retains the right leg $q$ and
its identity comparison.

The formal construction performs this lift on the universal family of
relative comparisons and evaluates at the given comparison. The resulting
identification carries its compatibility over $S$. This construction does
not require identifying it with an independently chosen pointwise lift.
\end{proof}

\begin{lemma}
\label[lemma]{lem:Supplementary_Reflecting_Base_Triangle}
Let $a,b\colon X\to D\times_T S$ be equipped with triangles
$\beta_a\colon q_Da=t$ and $\beta_b\colon q_Db=t$, where $q_D\colon D\times_T S\to S$ is the second projection and $t\colon X\to S$.
Suppose $\Phi$ compares their full induced pullback cones and its right
component comes with a specified identification
\[
 \beta_b\cdot\Phi_{\mathrm{right}}=\beta_a.
\]
Then $\Phi$ lifts to a comparison of $a$ and $b$ over $S$. The cone
comparison induced by this lift is identified with $\Phi$, including
its matching compatibility.
\end{lemma}
\begin{proof}[Proof idea]
Use the identification-anima universal property of the pullback to lift
$\Phi$ with its prescribed image. Its right projection computation transports
the displayed equation to the triangle equation of the lifted
identification. The full image computation also compares the matching
witnesses. Thus the same lift serves both the pullback calculation and the
calculation over the base.
\end{proof}

\begin{lemma}
\label[lemma]{lem:Supplementary_Base_Change_Laws}
There are comparisons over $S$
\[
 p^*(\operatorname{id})\simeq_S\operatorname{id},
 \qquad (p^*v)(p^*u)\simeq_S p^*(vu)
\]
for composable functors $u,v$ over $T$.
For composition, the resulting cone comparison recovers the comparison
obtained by pasting the defining pullback computations.
\end{lemma}
\begin{proof}[Proof idea]
For the identity, compare its defining cone with the original pullback
cone. The unitor on the first leg and the identity on the second satisfy
the matching equation by the triangle law.

For composition, restrict the defining cone of $p^*v$ along $p^*u$.
Apply its computation comparison, then the one for $p^*u$, and use the
restriction and composition comparisons for the action on cones. This
produces a full comparison with the defining cone of $p^*(vu)$.
Its right component satisfies the equation in
\Cref{lem:Supplementary_Reflecting_Base_Triangle}; checking that equation
uses the actual second-projection computations defining the triangles.
Reflection gives the desired comparison over $S$, together with its full
cone and triangle computations.
\end{proof}
\end{document}
